\section{Conclusion}
\label{sec:conclusion}
This study set out to determine whether fine-tuned RoBERTa-base detects deceptive hotel reviews more accurately than a fixed TF--IDF logistic-regression baseline when both models are evaluated on identical hotel-disjoint partitions, and at what cost.

\emph{RQ1.} The transformer did not significantly outperform the baseline. Its macro-F1 advantage on the primary partition was 0.0075 with $p=0.742$, and the hotel-grouped cross-validation gave a gain of 0.013 whose confidence interval includes zero. RoBERTa ranked reviews better, with higher average precision and ROC-AUC, but at the fixed threshold it traded missed deceptive reviews for false accusations against truthful ones. \emph{RQ2.} Neither model was consistently better across partitions or hotels: RoBERTa was slightly worse on two of the three partitions, won on 11 of 20 hotels, and the variation caused by changing the test hotels exceeded the difference between models. The models differed sharply in robustness, however. The baseline was unaffected by changes in case and whitespace, whereas RoBERTa lost up to 0.57 macro-F1 because those changes altered its tokenisation and pushed most reviews past its input limit. \emph{RQ3.} RoBERTa needed about 90 seconds of GPU training and classified about 133 reviews per second, compared with a fraction of a second of CPU training and more than 8,500 reviews per second for the baseline.

These findings matter in two ways. Methodologically, they show that the size of a transformer's advantage on a deception benchmark depends on how the data are partitioned and tested: with whole hotels held out and differences tested on paired predictions, the advantage is small, changes sign between partitions and is not statistically detectable. Reports of transformer superiority on small corpora should therefore be accompanied by a fixed baseline, grouped partitions and a paired test. Practically, a platform that screens large review streams could use the linear model as a fast, robust first filter, since it is two orders of magnitude cheaper and no less accurate in macro-F1 here. The transformer's better ranking would be useful only after threshold calibration and with input normalisation that protects it from formatting changes.

The study has clear limitations. It uses one small, balanced English corpus in which deceptive reviews were written by crowd workers on request; commercial fake reviews are produced differently, and the truthful labels are not independently verified. The corpus has no reviewer identifiers, so reviews by the same person may still appear in training and test data, and each fixed partition tests only three hotels, which limits the power of the prespecified test; the cross-validation and the review-level test that address this were added after the protocol. Hyperparameters and the threshold were fixed rather than tuned, RoBERTa truncates one in five reviews, and its training is not fully deterministic on the GPU. The planned YelpChi experiment and the star-rating ablation could not be run, so nothing is concluded about Yelp filtering decisions or the value of rating metadata.

Future work should follow directly from these limits. First, the same protocol should be applied to the text-containing YelpChi data, with reviewer-disjoint and chronological splits, to test whether the conclusions hold for commercially filtered reviews and to carry out the rating ablation. Second, calibrating the decision threshold on validation data would show whether RoBERTa's ranking advantage can become a macro-F1 advantage. Third, input normalisation, longer contexts and training with formatting augmentation should be evaluated against the perturbations used here, together with paraphrasing and machine-generated reviews. Finally, cross-domain tests on restaurant and medical reviews and cost measurements at realistic platform volumes would establish whether a lightweight baseline remains the better choice outside a single hotel benchmark.
